\scanpage{42}
EGA IV 15.6. Dans 15.6.1. (p.236 l. -18) supprimer \textquotedbl localement
noethérien\textquotedbl\ , cette hypothèse ne sert pas dans la démonstration; en revanche, dans 15.6.3. \inserted{et 15.6.4.} il faut al\deleted{\unsure{or|ro}}ors préciser qu'on suppose $f$ localement
de présentation finie, et dans \fbox{15.6.4. que de plus $f$ est quasi-séparé.}
% Hand-drawn frame, and oblique stroke between the first and third lines.

\handwritten{$\Bigg($}On notera que sous ces conditions, \deleted{la}le \deleted{proposition} corollaire 15.6.5.
devient inutile, étant déjà contenu dans ce qui précède.

IV 15.6.7. p.241, l. 2, après \textquotedbl no\deleted{\unsure{t}}ethérien\textquotedbl\ ajouter \textquotedbl ou si la condition
\handwritten{$\unsure{\mathrm{D}|1}$} de 15.6.6. est vérifiée\textquotedbl\ .

Il convient également d'ajouter quelques points d'\uline{autocritique}
(peut-être bloqués ensemble, en C) après les \textquotedbl modifications de textes),
savoir les points suivants:

1\textordmasculine) Suivant une idée de Raynaud, il apparait maintenant qu'on peut
traiter de façon nettement plus satisfaisante le co\deleted{\unsure{m}}ntenu des par.14 et
15 en utilisant l'outil de la hensélisation developpé au par.18 .
De cette façon, on arrive à des formulations et des démonstrations plus
simples, tous les résultats principaux étant de plus indépendants de toute hypothèse noethérienne. De plus, on arrive à établir aussi par cette
voie que la notion de morphisme universellement ouvert \deleted{\unsure{se compor}} de
présentation finie se comporte bien par limites projectives de schémas,
suivant le point de vue du par.8 . Nous espérons présenter une version
entièrement refondue de ces résultats dans la réédition, où \deleted{ils apparai\unsure{r|x}}
la substance des actuels par.14 et 15 se trouvera rejetée après le par.
18 (ce qui ne présente aucun inconvénient, les par. 16 et suivants
étant visiblement indépendants des actuels par.14 et 15)\unsure{/.|.}

2\textordmasculine) Nous nous répentons sincèrement de l'usage immodéré fait du
terme \textquotedbl non \uline{ramifié\textquotedbl\ au lieu de \textquotedbl net\textquotedbl} . Nous pourrions ajouter quelques
% archive sheet number: 41
